%% rset.conc.tex
%% Conclusions and future work (merged: detailed summary + the future-work
%% subsections formerly in rset.future).  Drop in via %% rset.conc.tex
%% Conclusions and future work (merged: detailed summary + the future-work
%% subsections formerly in rset.future).  Drop in via %% rset.conc.tex
%% Conclusions and future work (merged: detailed summary + the future-work
%% subsections formerly in rset.future).  Drop in via %% rset.conc.tex
%% Conclusions and future work (merged: detailed summary + the future-work
%% subsections formerly in rset.future).  Drop in via \input{rset.conc};
%% remove \input{rset.future} from rset.tex.
\input{rset.colors}   % red/black colour layer + FM symbols (providecommand-guarded)

\section{Conclusions and future work}
\label{sec:conc}
\label{sec:future}

We set out to reconcile two things usually held apart: a set theory with atoms,
of the kind nominal techniques require, and the austerity of building everything
from the empty set. The reconciliation is the knot. Taking two copies of set
theory --- a {\rd red} and a {\bk black} --- and tying them so that each colour's
atoms are the other colour's sets, we obtain an infinite supply of atoms without
positing a single one: every atom is a set of the dual colour, manufactured from
$\varnothing$, and the infinite risk that Chaitin's accounting charges to a
primitive infinity of atoms is refinanced as the ordinary risk of one empty set
and one successor.

Around this knot we have assembled a connected picture. The formal theory above
fixes the two copies and ties them; a self-contained sketch of
Fraenkel--Mostowski set theory (Section~\ref{sec:fm}) makes the target precise;
and the representation theorem (Section~\ref{sec:fmmodel}) exhibits the red
hereditarily finite sets over the black-set atoms as a model of the nominal
hereditarily finite sets, in which finite support --- imposed as an axiom in
$\mathsf{FM}$ --- is instead a theorem. The boundary between the finitary core and
its infinitary extension proved to be the exact ledger line between what the
construction \emph{derives} --- atoms, foundation, finite support --- and what it
must still \emph{assume}: a single $\omega$, and the finite-support cut, once
actual infinities are admitted.

The construction earns its keep in three further ways. It founds the
Pitts--Gabbay nominal techniques (Section~\ref{sec:applications}), and on exactly
the hereditarily finite fragment that the realised nominal languages inhabit, so
the restriction that bounds the set theory is the natural home of the computing;
it suggests, through reflective names, a more expressive theory of names in which
the nominal discipline reappears as an equivariant fragment; and it serves as a
bridge (Section~\ref{sec:nwf}) between the well-founded theories $\mathsf{ZF}$ and
$\mathsf{FM}$ and Aczel's non-well-founded universe, by making
non-well-foundedness not an axiom but an emergent, global property of two locally
well-founded theories. The road from this programmatic account to a rigorous one
runs, we have argued, through algebraic set theory (Section~\ref{sec:rigor}),
whose initial-$\mathsf{ZF}$-algebra theorem is purpose-built to supply the knotted
fixpoint and, with it, the membership and extensionality structure the
construction needs.

Several directions remain open; we record the ones we find most pressing.

\subsection{Arbitrarily many colours}
\label{sec:future:palette}

Two colours were a convenience, not a necessity. Replace the pair by any set $C$
of \emph{colours}; each $c\in C$ carries its own copy of the theory,
$\mathrm{Set}_c[X_c]$, with its own membership $\in_c$ and its own comprehension
$\{\cdots\}_c$. The only remaining datum is the \emph{nesting discipline}: which
colours may serve as atoms for which. This is a relation on $C$, and it is most
naturally presented as a directed graph --- equivalently a finite-state machine
--- with states $C$ and an edge $c\to c'$ whenever $c'$-sets are admissible
$c$-atoms. The atom supply of each colour is then the disjoint sum over its
outgoing edges,
\[
  X_c \;=\; \bigsqcup_{c\to c'}\ \mathrm{Set}_{c'}[X_{c'}],
\]
and the two-colour theory of this paper is the instance given by the two-state
automaton $\rd{\mathrm{red}} \rightleftarrows \bk{\mathrm{black}}$, with
$\rX=\bSetX$ and $\bX=\rSetX$.

Everything in the paper relativises to this picture. The well-formed cross-colour
nestings are the runs of the automaton; the local/global well-foundedness analysis
of Section~\ref{sec:nwf} reads off its structure, an acyclic colour graph giving a
globally well-founded universe and a cycle giving non-well-foundedness along that
cycle. (A self-loop $c\to c$ would let $c$-sets be their own atoms and collapse
the barrier; the two-cycle $\rd{\mathrm{red}}\rightleftarrows\bk{\mathrm{black}}$
is the smallest discipline that keeps every colour locally well-founded while
making the system globally non-well-founded.) This generalization is routine, and
we record it mainly to fix the slogan: \emph{a reflective set theory is an
automaton on colours}, and the construction above is its two-state instance.

\subsection{Generating colours dynamically}
\label{sec:future:dynamic}

The genuinely interesting question we leave open. A fixed automaton fixes the
palette in advance. One would instead like colours to be \emph{generated on
demand}: a fresh colour spawned in the course of building a set, the palette
growing as construction proceeds, the automaton unfolding rather than given. The
difficulty is not in admitting infinitely many colours --- a countable palette is
unremarkable --- but in doing so \emph{elegantly}, with a generation discipline
internal to the theory, so that the mechanism producing colours is of a piece with
the mechanism producing sets and names.

There is a suggestive circularity to exploit, which is why we think the question
worth posing rather than dismissing. Generating a fresh colour is generating a
fresh name, in exactly the sense that
Sections~\ref{sec:fm}--\ref{sec:applications} reconstruct: a colour is drawn from
a supply, distinct from those already in play, and otherwise structureless. So a
reflective theory might be made to \emph{generate its own colours} by the very
nominal machinery it models --- the palette drawn, like the atoms, from a stream,
and the colour automaton itself becoming a nominal object in a theory one level
up. Whether this reflexive step can be taken without infinite regress, what its
fixed point looks like, and whether it lands back inside the spectrum of
Section~\ref{sec:nwf}, are the questions we find most worth pursuing. We do not
settle them here; we mark them as the natural sequel to this work.

\subsection{A semantics for the rho calculus}
\label{sec:future:rho}

We close with a conjecture rather than a construction. Aczel gave a semantics for
Milner's $\mathsf{CCS}$ inside non-well-founded set theory \cite{aczel1988nwf}: a
process is identified with the set of its labelled transitions,
$\{\,(a,P') : P \xrightarrow{a} P'\,\}$, so that a looping process is a
non-well-founded set, bisimilar processes are equal, and the anti-foundation
axiom's solution lemma supplies and individuates processes as the unique
solutions of their behavioural equations. We surmise that the knotted universe
developed here is the natural home for the same construction carried out for the
reflective higher-order calculus, the rho calculus
\cite{DBLP:journals/entcs/MeredithR05}.

The fit is suggestive on two counts. First, the non-well-foundedness that Aczel
imports by axiom is here available for free, as the global, colour-blind structure
of the knotted universe (Section~\ref{sec:nwf}); the recursion that process
behaviour requires is the recursion the knot already carries. Second --- and this
is what sets rho apart from $\mathsf{CCS}$ and the $\pi$-calculus --- rho's names
are not primitive but \emph{reflected}: the quote operator $@$ turns a process
into a name, the dereference operator $*$ turns a name back into the process it
quotes, and $*@P \equiv P$. In the knotted universe a name may be taken to be an
\emph{atom}, and atoms here are constructed sets of the dual colour, opaque until
one crosses the colour barrier (Sections~\ref{sec:fmmodel:atoms}
and~\ref{sec:applications:reflective}). Reading $@$ as \emph{atom construction}
--- a process, presented as a set, is sealed into an atom of the opposite colour
--- and $*$ as \emph{atom unpacking} --- the barrier is crossed to recover the set
the atom seals --- the reflection law $*@P \equiv P$ becomes the round trip of
these two operations.

We do not carry the construction out, and we do not claim its details are routine:
communication, the substitution of names for names, and rho's equational theory
would each have to be matched against their set-theoretic shadows, and the
well-founded versus non-well-founded reading (Section~\ref{sec:rigor:fixpoint})
settled. We claim only that the ingredients line up --- a calculus whose names are
quoted processes asks for a set theory whose atoms are sets, and whose dynamics
asks for non-well-foundedness, and the knotted universe offers both at once --- and
we conjecture that such a semantics exists and is more convenient than one built
over an unstructured atom supply, precisely because the $@$/$*$ reflection it must
model is already the construction's defining move.
;
%% remove %% rset.future.tex
%% Future work: arbitrarily many colours, the finite-state-machine
%% generalization, and the open problem of generating colours dynamically.
%% Drop in via \input{rset.future}.
\input{rset.colors}   % red/black colour layer + FM symbols (providecommand-guarded)

\section{Future work: palettes and the generation of colour}
\label{sec:future}

\subsection{Arbitrarily many colours}
\label{sec:future:palette}

Two colours were a convenience, not a necessity. Replace the pair by any set $C$
of \emph{colours}; each $c\in C$ carries its own copy of the theory,
$\mathrm{Set}_c[X_c]$, with its own membership $\in_c$ and its own comprehension
$\{\cdots\}_c$. The only remaining datum is the \emph{nesting discipline}: which
colours may serve as atoms for which. This is a relation on $C$, and it is most
naturally presented as a directed graph --- equivalently a finite-state machine
--- with states $C$ and an edge $c\to c'$ whenever $c'$-sets are admissible
$c$-atoms. The atom supply of each colour is then the disjoint sum over its
outgoing edges,
\[
  X_c \;=\; \bigsqcup_{c\to c'}\ \mathrm{Set}_{c'}[X_{c'}],
\]
and the two-colour theory of this paper is the instance given by the two-state
automaton $\rd{\mathrm{red}} \rightleftarrows \bk{\mathrm{black}}$, with
$\rX=\bSetX$ and $\bX=\rSetX$.

Everything in the paper relativises to this picture. The well-formed cross-colour
nestings are the runs of the automaton; the local/global well-foundedness analysis
of Section~\ref{sec:nwf} reads off its structure, an acyclic colour graph giving a
globally well-founded universe and a cycle giving non-well-foundedness along that
cycle. (A self-loop $c\to c$ would let $c$-sets be their own atoms and collapse
the barrier; the two-cycle $\rd{\mathrm{red}}\rightleftarrows\bk{\mathrm{black}}$
is the smallest discipline that keeps every colour locally well-founded while
making the system globally non-well-founded.) This generalization is routine, and
we record it mainly to fix the slogan: \emph{a reflective set theory is an
automaton on colours}, and the constructions above are its two-state instance.

\subsection{Generating colours dynamically}
\label{sec:future:dynamic}

The genuinely interesting question we leave open. A fixed automaton fixes the
palette in advance. One would instead like colours to be \emph{generated on
demand}: a fresh colour spawned in the course of building a set, the palette
growing as construction proceeds, the automaton unfolding rather than given. The
difficulty is not in admitting infinitely many colours --- a countable palette is
unremarkable --- but in doing so \emph{elegantly}, with a generation discipline
internal to the theory, so that the mechanism producing colours is of a piece with
the mechanism producing sets and names.

There is a suggestive circularity to exploit, which is why we think the question
worth posing rather than dismissing. Generating a fresh colour is generating a
fresh name, in exactly the sense that
Sections~\ref{sec:fm}--\ref{sec:applications} reconstruct: a colour is drawn from
a supply, distinct from those already in play, and otherwise structureless. So a
reflective theory might be made to \emph{generate its own colours} by the very
nominal machinery it models --- the palette drawn, like the atoms, from a stream,
and the colour automaton itself becoming a nominal object in a theory one level
up. Whether this reflexive step can be taken without infinite regress, what its
fixed point looks like, and whether it lands back inside the spectrum of
Section~\ref{sec:nwf}, are the questions we find most worth pursuing. We do not
settle them here; we mark them as the natural sequel to this work.
 from rset.tex.
%% rset.colors.tex
%% Red/black colour layer for the two intertwined copies, plus the FM symbols
%% used by rset.mainthm.tex and rset.applications.tex.
%%
%% EVERYTHING routes through \rd (red copy = player) and \bk (black copy =
%% opponent): redefine just those two macros to change the whole scheme (e.g.
%% to make the "black" copy dark blue for legibility).  Every definition is
%% \providecommand-guarded, so it is safe to \input this file from several
%% places and harmless if rset.local already defines these names.
%%
%% Requires \usepackage{xcolor} (already loaded in rset.tex).

% --- the two copies -------------------------------------------------------
\providecommand{\rd}[1]{\textcolor{red}{#1}}     % red copy  (player)
\providecommand{\bk}[1]{\textcolor{black}{#1}}   % black copy (opponent)

% --- theory names and atom supplies --------------------------------------
\providecommand{\rSetX}{\rd{\mathrm{Set}[X]}}    % red   Set[X]
\providecommand{\bSetX}{\bk{\mathrm{Set}[X]}}    % black Set[X]
\providecommand{\rX}{\rd{X}}                      % red   atom supply
\providecommand{\bX}{\bk{X}}                      % black atom supply

% --- empty set ------------------------------------------------------------
\providecommand{\remp}{\rd{\emptyset}}
\providecommand{\bemp}{\bk{\emptyset}}

% --- membership (relations) ----------------------------------------------
\providecommand{\rin}{\mathrel{\rd{\in}}}
\providecommand{\bin}{\mathrel{\bk{\in}}}
\providecommand{\rnin}{\mathrel{\rd{\notin}}}
\providecommand{\bnin}{\mathrel{\bk{\notin}}}

% --- binary operations ----------------------------------------------------
\providecommand{\rcup}{\mathbin{\rd{\cup}}}
\providecommand{\bcup}{\mathbin{\bk{\cup}}}
\providecommand{\rcap}{\mathbin{\rd{\cap}}}
\providecommand{\bcap}{\mathbin{\bk{\cap}}}
\providecommand{\rsm}{\mathbin{\rd{\setminus}}}
\providecommand{\bsm}{\mathbin{\bk{\setminus}}}

% --- comprehension braces of each theory ---------------------------------
\providecommand{\rbr}[1]{\rd{\{}\,#1\,\rd{\}}}
\providecommand{\bbr}[1]{\bk{\{}\,#1\,\bk{\}}}

% --- colour-blind membership (sees through the colour barrier) ------------
\providecommand{\cbin}{\mathrel{\in_{\!\star}}}

% --- FM symbols (shared with rset.stdfmset.tex) --------------------------
\providecommand{\Atoms}{\mathbb{A}}
\providecommand{\Sym}{\mathrm{Sym}}
\providecommand{\pow}{\mathcal{P}}
\providecommand{\supp}{\mathrm{supp}}
\providecommand{\fs}{\mathrm{fs}}
\providecommand{\new}{\mathbin{\reflectbox{\ensuremath{\mathsf{N}}}}}
\providecommand{\atomsof}{\mathrm{atoms}}
\providecommand{\Red}{\rd{\mathsf{R}}}           % the red HF universe
   % red/black colour layer + FM symbols (providecommand-guarded)

\section{Conclusions and future work}
\label{sec:conc}
\label{sec:future}

We set out to reconcile two things usually held apart: a set theory with atoms,
of the kind nominal techniques require, and the austerity of building everything
from the empty set. The reconciliation is the knot. Taking two copies of set
theory --- a {\rd red} and a {\bk black} --- and tying them so that each colour's
atoms are the other colour's sets, we obtain an infinite supply of atoms without
positing a single one: every atom is a set of the dual colour, manufactured from
$\varnothing$, and the infinite risk that Chaitin's accounting charges to a
primitive infinity of atoms is refinanced as the ordinary risk of one empty set
and one successor.

Around this knot we have assembled a connected picture. The formal theory above
fixes the two copies and ties them; a self-contained sketch of
Fraenkel--Mostowski set theory (Section~\ref{sec:fm}) makes the target precise;
and the representation theorem (Section~\ref{sec:fmmodel}) exhibits the red
hereditarily finite sets over the black-set atoms as a model of the nominal
hereditarily finite sets, in which finite support --- imposed as an axiom in
$\mathsf{FM}$ --- is instead a theorem. The boundary between the finitary core and
its infinitary extension proved to be the exact ledger line between what the
construction \emph{derives} --- atoms, foundation, finite support --- and what it
must still \emph{assume}: a single $\omega$, and the finite-support cut, once
actual infinities are admitted.

The construction earns its keep in three further ways. It founds the
Pitts--Gabbay nominal techniques (Section~\ref{sec:applications}), and on exactly
the hereditarily finite fragment that the realised nominal languages inhabit, so
the restriction that bounds the set theory is the natural home of the computing;
it suggests, through reflective names, a more expressive theory of names in which
the nominal discipline reappears as an equivariant fragment; and it serves as a
bridge (Section~\ref{sec:nwf}) between the well-founded theories $\mathsf{ZF}$ and
$\mathsf{FM}$ and Aczel's non-well-founded universe, by making
non-well-foundedness not an axiom but an emergent, global property of two locally
well-founded theories. The road from this programmatic account to a rigorous one
runs, we have argued, through algebraic set theory (Section~\ref{sec:rigor}),
whose initial-$\mathsf{ZF}$-algebra theorem is purpose-built to supply the knotted
fixpoint and, with it, the membership and extensionality structure the
construction needs.

Several directions remain open; we record the ones we find most pressing.

\subsection{Arbitrarily many colours}
\label{sec:future:palette}

Two colours were a convenience, not a necessity. Replace the pair by any set $C$
of \emph{colours}; each $c\in C$ carries its own copy of the theory,
$\mathrm{Set}_c[X_c]$, with its own membership $\in_c$ and its own comprehension
$\{\cdots\}_c$. The only remaining datum is the \emph{nesting discipline}: which
colours may serve as atoms for which. This is a relation on $C$, and it is most
naturally presented as a directed graph --- equivalently a finite-state machine
--- with states $C$ and an edge $c\to c'$ whenever $c'$-sets are admissible
$c$-atoms. The atom supply of each colour is then the disjoint sum over its
outgoing edges,
\[
  X_c \;=\; \bigsqcup_{c\to c'}\ \mathrm{Set}_{c'}[X_{c'}],
\]
and the two-colour theory of this paper is the instance given by the two-state
automaton $\rd{\mathrm{red}} \rightleftarrows \bk{\mathrm{black}}$, with
$\rX=\bSetX$ and $\bX=\rSetX$.

Everything in the paper relativises to this picture. The well-formed cross-colour
nestings are the runs of the automaton; the local/global well-foundedness analysis
of Section~\ref{sec:nwf} reads off its structure, an acyclic colour graph giving a
globally well-founded universe and a cycle giving non-well-foundedness along that
cycle. (A self-loop $c\to c$ would let $c$-sets be their own atoms and collapse
the barrier; the two-cycle $\rd{\mathrm{red}}\rightleftarrows\bk{\mathrm{black}}$
is the smallest discipline that keeps every colour locally well-founded while
making the system globally non-well-founded.) This generalization is routine, and
we record it mainly to fix the slogan: \emph{a reflective set theory is an
automaton on colours}, and the construction above is its two-state instance.

\subsection{Generating colours dynamically}
\label{sec:future:dynamic}

The genuinely interesting question we leave open. A fixed automaton fixes the
palette in advance. One would instead like colours to be \emph{generated on
demand}: a fresh colour spawned in the course of building a set, the palette
growing as construction proceeds, the automaton unfolding rather than given. The
difficulty is not in admitting infinitely many colours --- a countable palette is
unremarkable --- but in doing so \emph{elegantly}, with a generation discipline
internal to the theory, so that the mechanism producing colours is of a piece with
the mechanism producing sets and names.

There is a suggestive circularity to exploit, which is why we think the question
worth posing rather than dismissing. Generating a fresh colour is generating a
fresh name, in exactly the sense that
Sections~\ref{sec:fm}--\ref{sec:applications} reconstruct: a colour is drawn from
a supply, distinct from those already in play, and otherwise structureless. So a
reflective theory might be made to \emph{generate its own colours} by the very
nominal machinery it models --- the palette drawn, like the atoms, from a stream,
and the colour automaton itself becoming a nominal object in a theory one level
up. Whether this reflexive step can be taken without infinite regress, what its
fixed point looks like, and whether it lands back inside the spectrum of
Section~\ref{sec:nwf}, are the questions we find most worth pursuing. We do not
settle them here; we mark them as the natural sequel to this work.

\subsection{A semantics for the rho calculus}
\label{sec:future:rho}

We close with a conjecture rather than a construction. Aczel gave a semantics for
Milner's $\mathsf{CCS}$ inside non-well-founded set theory \cite{aczel1988nwf}: a
process is identified with the set of its labelled transitions,
$\{\,(a,P') : P \xrightarrow{a} P'\,\}$, so that a looping process is a
non-well-founded set, bisimilar processes are equal, and the anti-foundation
axiom's solution lemma supplies and individuates processes as the unique
solutions of their behavioural equations. We surmise that the knotted universe
developed here is the natural home for the same construction carried out for the
reflective higher-order calculus, the rho calculus
\cite{DBLP:journals/entcs/MeredithR05}.

The fit is suggestive on two counts. First, the non-well-foundedness that Aczel
imports by axiom is here available for free, as the global, colour-blind structure
of the knotted universe (Section~\ref{sec:nwf}); the recursion that process
behaviour requires is the recursion the knot already carries. Second --- and this
is what sets rho apart from $\mathsf{CCS}$ and the $\pi$-calculus --- rho's names
are not primitive but \emph{reflected}: the quote operator $@$ turns a process
into a name, the dereference operator $*$ turns a name back into the process it
quotes, and $*@P \equiv P$. In the knotted universe a name may be taken to be an
\emph{atom}, and atoms here are constructed sets of the dual colour, opaque until
one crosses the colour barrier (Sections~\ref{sec:fmmodel:atoms}
and~\ref{sec:applications:reflective}). Reading $@$ as \emph{atom construction}
--- a process, presented as a set, is sealed into an atom of the opposite colour
--- and $*$ as \emph{atom unpacking} --- the barrier is crossed to recover the set
the atom seals --- the reflection law $*@P \equiv P$ becomes the round trip of
these two operations.

We do not carry the construction out, and we do not claim its details are routine:
communication, the substitution of names for names, and rho's equational theory
would each have to be matched against their set-theoretic shadows, and the
well-founded versus non-well-founded reading (Section~\ref{sec:rigor:fixpoint})
settled. We claim only that the ingredients line up --- a calculus whose names are
quoted processes asks for a set theory whose atoms are sets, and whose dynamics
asks for non-well-foundedness, and the knotted universe offers both at once --- and
we conjecture that such a semantics exists and is more convenient than one built
over an unstructured atom supply, precisely because the $@$/$*$ reflection it must
model is already the construction's defining move.
;
%% remove %% rset.future.tex
%% Future work: arbitrarily many colours, the finite-state-machine
%% generalization, and the open problem of generating colours dynamically.
%% Drop in via %% rset.future.tex
%% Future work: arbitrarily many colours, the finite-state-machine
%% generalization, and the open problem of generating colours dynamically.
%% Drop in via \input{rset.future}.
\input{rset.colors}   % red/black colour layer + FM symbols (providecommand-guarded)

\section{Future work: palettes and the generation of colour}
\label{sec:future}

\subsection{Arbitrarily many colours}
\label{sec:future:palette}

Two colours were a convenience, not a necessity. Replace the pair by any set $C$
of \emph{colours}; each $c\in C$ carries its own copy of the theory,
$\mathrm{Set}_c[X_c]$, with its own membership $\in_c$ and its own comprehension
$\{\cdots\}_c$. The only remaining datum is the \emph{nesting discipline}: which
colours may serve as atoms for which. This is a relation on $C$, and it is most
naturally presented as a directed graph --- equivalently a finite-state machine
--- with states $C$ and an edge $c\to c'$ whenever $c'$-sets are admissible
$c$-atoms. The atom supply of each colour is then the disjoint sum over its
outgoing edges,
\[
  X_c \;=\; \bigsqcup_{c\to c'}\ \mathrm{Set}_{c'}[X_{c'}],
\]
and the two-colour theory of this paper is the instance given by the two-state
automaton $\rd{\mathrm{red}} \rightleftarrows \bk{\mathrm{black}}$, with
$\rX=\bSetX$ and $\bX=\rSetX$.

Everything in the paper relativises to this picture. The well-formed cross-colour
nestings are the runs of the automaton; the local/global well-foundedness analysis
of Section~\ref{sec:nwf} reads off its structure, an acyclic colour graph giving a
globally well-founded universe and a cycle giving non-well-foundedness along that
cycle. (A self-loop $c\to c$ would let $c$-sets be their own atoms and collapse
the barrier; the two-cycle $\rd{\mathrm{red}}\rightleftarrows\bk{\mathrm{black}}$
is the smallest discipline that keeps every colour locally well-founded while
making the system globally non-well-founded.) This generalization is routine, and
we record it mainly to fix the slogan: \emph{a reflective set theory is an
automaton on colours}, and the constructions above are its two-state instance.

\subsection{Generating colours dynamically}
\label{sec:future:dynamic}

The genuinely interesting question we leave open. A fixed automaton fixes the
palette in advance. One would instead like colours to be \emph{generated on
demand}: a fresh colour spawned in the course of building a set, the palette
growing as construction proceeds, the automaton unfolding rather than given. The
difficulty is not in admitting infinitely many colours --- a countable palette is
unremarkable --- but in doing so \emph{elegantly}, with a generation discipline
internal to the theory, so that the mechanism producing colours is of a piece with
the mechanism producing sets and names.

There is a suggestive circularity to exploit, which is why we think the question
worth posing rather than dismissing. Generating a fresh colour is generating a
fresh name, in exactly the sense that
Sections~\ref{sec:fm}--\ref{sec:applications} reconstruct: a colour is drawn from
a supply, distinct from those already in play, and otherwise structureless. So a
reflective theory might be made to \emph{generate its own colours} by the very
nominal machinery it models --- the palette drawn, like the atoms, from a stream,
and the colour automaton itself becoming a nominal object in a theory one level
up. Whether this reflexive step can be taken without infinite regress, what its
fixed point looks like, and whether it lands back inside the spectrum of
Section~\ref{sec:nwf}, are the questions we find most worth pursuing. We do not
settle them here; we mark them as the natural sequel to this work.
.
%% rset.colors.tex
%% Red/black colour layer for the two intertwined copies, plus the FM symbols
%% used by rset.mainthm.tex and rset.applications.tex.
%%
%% EVERYTHING routes through \rd (red copy = player) and \bk (black copy =
%% opponent): redefine just those two macros to change the whole scheme (e.g.
%% to make the "black" copy dark blue for legibility).  Every definition is
%% \providecommand-guarded, so it is safe to \input this file from several
%% places and harmless if rset.local already defines these names.
%%
%% Requires \usepackage{xcolor} (already loaded in rset.tex).

% --- the two copies -------------------------------------------------------
\providecommand{\rd}[1]{\textcolor{red}{#1}}     % red copy  (player)
\providecommand{\bk}[1]{\textcolor{black}{#1}}   % black copy (opponent)

% --- theory names and atom supplies --------------------------------------
\providecommand{\rSetX}{\rd{\mathrm{Set}[X]}}    % red   Set[X]
\providecommand{\bSetX}{\bk{\mathrm{Set}[X]}}    % black Set[X]
\providecommand{\rX}{\rd{X}}                      % red   atom supply
\providecommand{\bX}{\bk{X}}                      % black atom supply

% --- empty set ------------------------------------------------------------
\providecommand{\remp}{\rd{\emptyset}}
\providecommand{\bemp}{\bk{\emptyset}}

% --- membership (relations) ----------------------------------------------
\providecommand{\rin}{\mathrel{\rd{\in}}}
\providecommand{\bin}{\mathrel{\bk{\in}}}
\providecommand{\rnin}{\mathrel{\rd{\notin}}}
\providecommand{\bnin}{\mathrel{\bk{\notin}}}

% --- binary operations ----------------------------------------------------
\providecommand{\rcup}{\mathbin{\rd{\cup}}}
\providecommand{\bcup}{\mathbin{\bk{\cup}}}
\providecommand{\rcap}{\mathbin{\rd{\cap}}}
\providecommand{\bcap}{\mathbin{\bk{\cap}}}
\providecommand{\rsm}{\mathbin{\rd{\setminus}}}
\providecommand{\bsm}{\mathbin{\bk{\setminus}}}

% --- comprehension braces of each theory ---------------------------------
\providecommand{\rbr}[1]{\rd{\{}\,#1\,\rd{\}}}
\providecommand{\bbr}[1]{\bk{\{}\,#1\,\bk{\}}}

% --- colour-blind membership (sees through the colour barrier) ------------
\providecommand{\cbin}{\mathrel{\in_{\!\star}}}

% --- FM symbols (shared with rset.stdfmset.tex) --------------------------
\providecommand{\Atoms}{\mathbb{A}}
\providecommand{\Sym}{\mathrm{Sym}}
\providecommand{\pow}{\mathcal{P}}
\providecommand{\supp}{\mathrm{supp}}
\providecommand{\fs}{\mathrm{fs}}
\providecommand{\new}{\mathbin{\reflectbox{\ensuremath{\mathsf{N}}}}}
\providecommand{\atomsof}{\mathrm{atoms}}
\providecommand{\Red}{\rd{\mathsf{R}}}           % the red HF universe
   % red/black colour layer + FM symbols (providecommand-guarded)

\section{Future work: palettes and the generation of colour}
\label{sec:future}

\subsection{Arbitrarily many colours}
\label{sec:future:palette}

Two colours were a convenience, not a necessity. Replace the pair by any set $C$
of \emph{colours}; each $c\in C$ carries its own copy of the theory,
$\mathrm{Set}_c[X_c]$, with its own membership $\in_c$ and its own comprehension
$\{\cdots\}_c$. The only remaining datum is the \emph{nesting discipline}: which
colours may serve as atoms for which. This is a relation on $C$, and it is most
naturally presented as a directed graph --- equivalently a finite-state machine
--- with states $C$ and an edge $c\to c'$ whenever $c'$-sets are admissible
$c$-atoms. The atom supply of each colour is then the disjoint sum over its
outgoing edges,
\[
  X_c \;=\; \bigsqcup_{c\to c'}\ \mathrm{Set}_{c'}[X_{c'}],
\]
and the two-colour theory of this paper is the instance given by the two-state
automaton $\rd{\mathrm{red}} \rightleftarrows \bk{\mathrm{black}}$, with
$\rX=\bSetX$ and $\bX=\rSetX$.

Everything in the paper relativises to this picture. The well-formed cross-colour
nestings are the runs of the automaton; the local/global well-foundedness analysis
of Section~\ref{sec:nwf} reads off its structure, an acyclic colour graph giving a
globally well-founded universe and a cycle giving non-well-foundedness along that
cycle. (A self-loop $c\to c$ would let $c$-sets be their own atoms and collapse
the barrier; the two-cycle $\rd{\mathrm{red}}\rightleftarrows\bk{\mathrm{black}}$
is the smallest discipline that keeps every colour locally well-founded while
making the system globally non-well-founded.) This generalization is routine, and
we record it mainly to fix the slogan: \emph{a reflective set theory is an
automaton on colours}, and the constructions above are its two-state instance.

\subsection{Generating colours dynamically}
\label{sec:future:dynamic}

The genuinely interesting question we leave open. A fixed automaton fixes the
palette in advance. One would instead like colours to be \emph{generated on
demand}: a fresh colour spawned in the course of building a set, the palette
growing as construction proceeds, the automaton unfolding rather than given. The
difficulty is not in admitting infinitely many colours --- a countable palette is
unremarkable --- but in doing so \emph{elegantly}, with a generation discipline
internal to the theory, so that the mechanism producing colours is of a piece with
the mechanism producing sets and names.

There is a suggestive circularity to exploit, which is why we think the question
worth posing rather than dismissing. Generating a fresh colour is generating a
fresh name, in exactly the sense that
Sections~\ref{sec:fm}--\ref{sec:applications} reconstruct: a colour is drawn from
a supply, distinct from those already in play, and otherwise structureless. So a
reflective theory might be made to \emph{generate its own colours} by the very
nominal machinery it models --- the palette drawn, like the atoms, from a stream,
and the colour automaton itself becoming a nominal object in a theory one level
up. Whether this reflexive step can be taken without infinite regress, what its
fixed point looks like, and whether it lands back inside the spectrum of
Section~\ref{sec:nwf}, are the questions we find most worth pursuing. We do not
settle them here; we mark them as the natural sequel to this work.
 from rset.tex.
%% rset.colors.tex
%% Red/black colour layer for the two intertwined copies, plus the FM symbols
%% used by rset.mainthm.tex and rset.applications.tex.
%%
%% EVERYTHING routes through \rd (red copy = player) and \bk (black copy =
%% opponent): redefine just those two macros to change the whole scheme (e.g.
%% to make the "black" copy dark blue for legibility).  Every definition is
%% \providecommand-guarded, so it is safe to \input this file from several
%% places and harmless if rset.local already defines these names.
%%
%% Requires \usepackage{xcolor} (already loaded in rset.tex).

% --- the two copies -------------------------------------------------------
\providecommand{\rd}[1]{\textcolor{red}{#1}}     % red copy  (player)
\providecommand{\bk}[1]{\textcolor{black}{#1}}   % black copy (opponent)

% --- theory names and atom supplies --------------------------------------
\providecommand{\rSetX}{\rd{\mathrm{Set}[X]}}    % red   Set[X]
\providecommand{\bSetX}{\bk{\mathrm{Set}[X]}}    % black Set[X]
\providecommand{\rX}{\rd{X}}                      % red   atom supply
\providecommand{\bX}{\bk{X}}                      % black atom supply

% --- empty set ------------------------------------------------------------
\providecommand{\remp}{\rd{\emptyset}}
\providecommand{\bemp}{\bk{\emptyset}}

% --- membership (relations) ----------------------------------------------
\providecommand{\rin}{\mathrel{\rd{\in}}}
\providecommand{\bin}{\mathrel{\bk{\in}}}
\providecommand{\rnin}{\mathrel{\rd{\notin}}}
\providecommand{\bnin}{\mathrel{\bk{\notin}}}

% --- binary operations ----------------------------------------------------
\providecommand{\rcup}{\mathbin{\rd{\cup}}}
\providecommand{\bcup}{\mathbin{\bk{\cup}}}
\providecommand{\rcap}{\mathbin{\rd{\cap}}}
\providecommand{\bcap}{\mathbin{\bk{\cap}}}
\providecommand{\rsm}{\mathbin{\rd{\setminus}}}
\providecommand{\bsm}{\mathbin{\bk{\setminus}}}

% --- comprehension braces of each theory ---------------------------------
\providecommand{\rbr}[1]{\rd{\{}\,#1\,\rd{\}}}
\providecommand{\bbr}[1]{\bk{\{}\,#1\,\bk{\}}}

% --- colour-blind membership (sees through the colour barrier) ------------
\providecommand{\cbin}{\mathrel{\in_{\!\star}}}

% --- FM symbols (shared with rset.stdfmset.tex) --------------------------
\providecommand{\Atoms}{\mathbb{A}}
\providecommand{\Sym}{\mathrm{Sym}}
\providecommand{\pow}{\mathcal{P}}
\providecommand{\supp}{\mathrm{supp}}
\providecommand{\fs}{\mathrm{fs}}
\providecommand{\new}{\mathbin{\reflectbox{\ensuremath{\mathsf{N}}}}}
\providecommand{\atomsof}{\mathrm{atoms}}
\providecommand{\Red}{\rd{\mathsf{R}}}           % the red HF universe
   % red/black colour layer + FM symbols (providecommand-guarded)

\section{Conclusions and future work}
\label{sec:conc}
\label{sec:future}

We set out to reconcile two things usually held apart: a set theory with atoms,
of the kind nominal techniques require, and the austerity of building everything
from the empty set. The reconciliation is the knot. Taking two copies of set
theory --- a {\rd red} and a {\bk black} --- and tying them so that each colour's
atoms are the other colour's sets, we obtain an infinite supply of atoms without
positing a single one: every atom is a set of the dual colour, manufactured from
$\varnothing$, and the infinite risk that Chaitin's accounting charges to a
primitive infinity of atoms is refinanced as the ordinary risk of one empty set
and one successor.

Around this knot we have assembled a connected picture. The formal theory above
fixes the two copies and ties them; a self-contained sketch of
Fraenkel--Mostowski set theory (Section~\ref{sec:fm}) makes the target precise;
and the representation theorem (Section~\ref{sec:fmmodel}) exhibits the red
hereditarily finite sets over the black-set atoms as a model of the nominal
hereditarily finite sets, in which finite support --- imposed as an axiom in
$\mathsf{FM}$ --- is instead a theorem. The boundary between the finitary core and
its infinitary extension proved to be the exact ledger line between what the
construction \emph{derives} --- atoms, foundation, finite support --- and what it
must still \emph{assume}: a single $\omega$, and the finite-support cut, once
actual infinities are admitted.

The construction earns its keep in three further ways. It founds the
Pitts--Gabbay nominal techniques (Section~\ref{sec:applications}), and on exactly
the hereditarily finite fragment that the realised nominal languages inhabit, so
the restriction that bounds the set theory is the natural home of the computing;
it suggests, through reflective names, a more expressive theory of names in which
the nominal discipline reappears as an equivariant fragment; and it serves as a
bridge (Section~\ref{sec:nwf}) between the well-founded theories $\mathsf{ZF}$ and
$\mathsf{FM}$ and Aczel's non-well-founded universe, by making
non-well-foundedness not an axiom but an emergent, global property of two locally
well-founded theories. The road from this programmatic account to a rigorous one
runs, we have argued, through algebraic set theory (Section~\ref{sec:rigor}),
whose initial-$\mathsf{ZF}$-algebra theorem is purpose-built to supply the knotted
fixpoint and, with it, the membership and extensionality structure the
construction needs.

Several directions remain open; we record the ones we find most pressing.

\subsection{Arbitrarily many colours}
\label{sec:future:palette}

Two colours were a convenience, not a necessity. Replace the pair by any set $C$
of \emph{colours}; each $c\in C$ carries its own copy of the theory,
$\mathrm{Set}_c[X_c]$, with its own membership $\in_c$ and its own comprehension
$\{\cdots\}_c$. The only remaining datum is the \emph{nesting discipline}: which
colours may serve as atoms for which. This is a relation on $C$, and it is most
naturally presented as a directed graph --- equivalently a finite-state machine
--- with states $C$ and an edge $c\to c'$ whenever $c'$-sets are admissible
$c$-atoms. The atom supply of each colour is then the disjoint sum over its
outgoing edges,
\[
  X_c \;=\; \bigsqcup_{c\to c'}\ \mathrm{Set}_{c'}[X_{c'}],
\]
and the two-colour theory of this paper is the instance given by the two-state
automaton $\rd{\mathrm{red}} \rightleftarrows \bk{\mathrm{black}}$, with
$\rX=\bSetX$ and $\bX=\rSetX$.

Everything in the paper relativises to this picture. The well-formed cross-colour
nestings are the runs of the automaton; the local/global well-foundedness analysis
of Section~\ref{sec:nwf} reads off its structure, an acyclic colour graph giving a
globally well-founded universe and a cycle giving non-well-foundedness along that
cycle. (A self-loop $c\to c$ would let $c$-sets be their own atoms and collapse
the barrier; the two-cycle $\rd{\mathrm{red}}\rightleftarrows\bk{\mathrm{black}}$
is the smallest discipline that keeps every colour locally well-founded while
making the system globally non-well-founded.) This generalization is routine, and
we record it mainly to fix the slogan: \emph{a reflective set theory is an
automaton on colours}, and the construction above is its two-state instance.

\subsection{Generating colours dynamically}
\label{sec:future:dynamic}

The genuinely interesting question we leave open. A fixed automaton fixes the
palette in advance. One would instead like colours to be \emph{generated on
demand}: a fresh colour spawned in the course of building a set, the palette
growing as construction proceeds, the automaton unfolding rather than given. The
difficulty is not in admitting infinitely many colours --- a countable palette is
unremarkable --- but in doing so \emph{elegantly}, with a generation discipline
internal to the theory, so that the mechanism producing colours is of a piece with
the mechanism producing sets and names.

There is a suggestive circularity to exploit, which is why we think the question
worth posing rather than dismissing. Generating a fresh colour is generating a
fresh name, in exactly the sense that
Sections~\ref{sec:fm}--\ref{sec:applications} reconstruct: a colour is drawn from
a supply, distinct from those already in play, and otherwise structureless. So a
reflective theory might be made to \emph{generate its own colours} by the very
nominal machinery it models --- the palette drawn, like the atoms, from a stream,
and the colour automaton itself becoming a nominal object in a theory one level
up. Whether this reflexive step can be taken without infinite regress, what its
fixed point looks like, and whether it lands back inside the spectrum of
Section~\ref{sec:nwf}, are the questions we find most worth pursuing. We do not
settle them here; we mark them as the natural sequel to this work.

\subsection{A semantics for the rho calculus}
\label{sec:future:rho}

We close with a conjecture rather than a construction. Aczel gave a semantics for
Milner's $\mathsf{CCS}$ inside non-well-founded set theory \cite{aczel1988nwf}: a
process is identified with the set of its labelled transitions,
$\{\,(a,P') : P \xrightarrow{a} P'\,\}$, so that a looping process is a
non-well-founded set, bisimilar processes are equal, and the anti-foundation
axiom's solution lemma supplies and individuates processes as the unique
solutions of their behavioural equations. We surmise that the knotted universe
developed here is the natural home for the same construction carried out for the
reflective higher-order calculus, the rho calculus
\cite{DBLP:journals/entcs/MeredithR05}.

The fit is suggestive on two counts. First, the non-well-foundedness that Aczel
imports by axiom is here available for free, as the global, colour-blind structure
of the knotted universe (Section~\ref{sec:nwf}); the recursion that process
behaviour requires is the recursion the knot already carries. Second --- and this
is what sets rho apart from $\mathsf{CCS}$ and the $\pi$-calculus --- rho's names
are not primitive but \emph{reflected}: the quote operator $@$ turns a process
into a name, the dereference operator $*$ turns a name back into the process it
quotes, and $*@P \equiv P$. In the knotted universe a name may be taken to be an
\emph{atom}, and atoms here are constructed sets of the dual colour, opaque until
one crosses the colour barrier (Sections~\ref{sec:fmmodel:atoms}
and~\ref{sec:applications:reflective}). Reading $@$ as \emph{atom construction}
--- a process, presented as a set, is sealed into an atom of the opposite colour
--- and $*$ as \emph{atom unpacking} --- the barrier is crossed to recover the set
the atom seals --- the reflection law $*@P \equiv P$ becomes the round trip of
these two operations.

We do not carry the construction out, and we do not claim its details are routine:
communication, the substitution of names for names, and rho's equational theory
would each have to be matched against their set-theoretic shadows, and the
well-founded versus non-well-founded reading (Section~\ref{sec:rigor:fixpoint})
settled. We claim only that the ingredients line up --- a calculus whose names are
quoted processes asks for a set theory whose atoms are sets, and whose dynamics
asks for non-well-foundedness, and the knotted universe offers both at once --- and
we conjecture that such a semantics exists and is more convenient than one built
over an unstructured atom supply, precisely because the $@$/$*$ reflection it must
model is already the construction's defining move.
;
%% remove %% rset.future.tex
%% Future work: arbitrarily many colours, the finite-state-machine
%% generalization, and the open problem of generating colours dynamically.
%% Drop in via %% rset.future.tex
%% Future work: arbitrarily many colours, the finite-state-machine
%% generalization, and the open problem of generating colours dynamically.
%% Drop in via %% rset.future.tex
%% Future work: arbitrarily many colours, the finite-state-machine
%% generalization, and the open problem of generating colours dynamically.
%% Drop in via \input{rset.future}.
\input{rset.colors}   % red/black colour layer + FM symbols (providecommand-guarded)

\section{Future work: palettes and the generation of colour}
\label{sec:future}

\subsection{Arbitrarily many colours}
\label{sec:future:palette}

Two colours were a convenience, not a necessity. Replace the pair by any set $C$
of \emph{colours}; each $c\in C$ carries its own copy of the theory,
$\mathrm{Set}_c[X_c]$, with its own membership $\in_c$ and its own comprehension
$\{\cdots\}_c$. The only remaining datum is the \emph{nesting discipline}: which
colours may serve as atoms for which. This is a relation on $C$, and it is most
naturally presented as a directed graph --- equivalently a finite-state machine
--- with states $C$ and an edge $c\to c'$ whenever $c'$-sets are admissible
$c$-atoms. The atom supply of each colour is then the disjoint sum over its
outgoing edges,
\[
  X_c \;=\; \bigsqcup_{c\to c'}\ \mathrm{Set}_{c'}[X_{c'}],
\]
and the two-colour theory of this paper is the instance given by the two-state
automaton $\rd{\mathrm{red}} \rightleftarrows \bk{\mathrm{black}}$, with
$\rX=\bSetX$ and $\bX=\rSetX$.

Everything in the paper relativises to this picture. The well-formed cross-colour
nestings are the runs of the automaton; the local/global well-foundedness analysis
of Section~\ref{sec:nwf} reads off its structure, an acyclic colour graph giving a
globally well-founded universe and a cycle giving non-well-foundedness along that
cycle. (A self-loop $c\to c$ would let $c$-sets be their own atoms and collapse
the barrier; the two-cycle $\rd{\mathrm{red}}\rightleftarrows\bk{\mathrm{black}}$
is the smallest discipline that keeps every colour locally well-founded while
making the system globally non-well-founded.) This generalization is routine, and
we record it mainly to fix the slogan: \emph{a reflective set theory is an
automaton on colours}, and the constructions above are its two-state instance.

\subsection{Generating colours dynamically}
\label{sec:future:dynamic}

The genuinely interesting question we leave open. A fixed automaton fixes the
palette in advance. One would instead like colours to be \emph{generated on
demand}: a fresh colour spawned in the course of building a set, the palette
growing as construction proceeds, the automaton unfolding rather than given. The
difficulty is not in admitting infinitely many colours --- a countable palette is
unremarkable --- but in doing so \emph{elegantly}, with a generation discipline
internal to the theory, so that the mechanism producing colours is of a piece with
the mechanism producing sets and names.

There is a suggestive circularity to exploit, which is why we think the question
worth posing rather than dismissing. Generating a fresh colour is generating a
fresh name, in exactly the sense that
Sections~\ref{sec:fm}--\ref{sec:applications} reconstruct: a colour is drawn from
a supply, distinct from those already in play, and otherwise structureless. So a
reflective theory might be made to \emph{generate its own colours} by the very
nominal machinery it models --- the palette drawn, like the atoms, from a stream,
and the colour automaton itself becoming a nominal object in a theory one level
up. Whether this reflexive step can be taken without infinite regress, what its
fixed point looks like, and whether it lands back inside the spectrum of
Section~\ref{sec:nwf}, are the questions we find most worth pursuing. We do not
settle them here; we mark them as the natural sequel to this work.
.
%% rset.colors.tex
%% Red/black colour layer for the two intertwined copies, plus the FM symbols
%% used by rset.mainthm.tex and rset.applications.tex.
%%
%% EVERYTHING routes through \rd (red copy = player) and \bk (black copy =
%% opponent): redefine just those two macros to change the whole scheme (e.g.
%% to make the "black" copy dark blue for legibility).  Every definition is
%% \providecommand-guarded, so it is safe to \input this file from several
%% places and harmless if rset.local already defines these names.
%%
%% Requires \usepackage{xcolor} (already loaded in rset.tex).

% --- the two copies -------------------------------------------------------
\providecommand{\rd}[1]{\textcolor{red}{#1}}     % red copy  (player)
\providecommand{\bk}[1]{\textcolor{black}{#1}}   % black copy (opponent)

% --- theory names and atom supplies --------------------------------------
\providecommand{\rSetX}{\rd{\mathrm{Set}[X]}}    % red   Set[X]
\providecommand{\bSetX}{\bk{\mathrm{Set}[X]}}    % black Set[X]
\providecommand{\rX}{\rd{X}}                      % red   atom supply
\providecommand{\bX}{\bk{X}}                      % black atom supply

% --- empty set ------------------------------------------------------------
\providecommand{\remp}{\rd{\emptyset}}
\providecommand{\bemp}{\bk{\emptyset}}

% --- membership (relations) ----------------------------------------------
\providecommand{\rin}{\mathrel{\rd{\in}}}
\providecommand{\bin}{\mathrel{\bk{\in}}}
\providecommand{\rnin}{\mathrel{\rd{\notin}}}
\providecommand{\bnin}{\mathrel{\bk{\notin}}}

% --- binary operations ----------------------------------------------------
\providecommand{\rcup}{\mathbin{\rd{\cup}}}
\providecommand{\bcup}{\mathbin{\bk{\cup}}}
\providecommand{\rcap}{\mathbin{\rd{\cap}}}
\providecommand{\bcap}{\mathbin{\bk{\cap}}}
\providecommand{\rsm}{\mathbin{\rd{\setminus}}}
\providecommand{\bsm}{\mathbin{\bk{\setminus}}}

% --- comprehension braces of each theory ---------------------------------
\providecommand{\rbr}[1]{\rd{\{}\,#1\,\rd{\}}}
\providecommand{\bbr}[1]{\bk{\{}\,#1\,\bk{\}}}

% --- colour-blind membership (sees through the colour barrier) ------------
\providecommand{\cbin}{\mathrel{\in_{\!\star}}}

% --- FM symbols (shared with rset.stdfmset.tex) --------------------------
\providecommand{\Atoms}{\mathbb{A}}
\providecommand{\Sym}{\mathrm{Sym}}
\providecommand{\pow}{\mathcal{P}}
\providecommand{\supp}{\mathrm{supp}}
\providecommand{\fs}{\mathrm{fs}}
\providecommand{\new}{\mathbin{\reflectbox{\ensuremath{\mathsf{N}}}}}
\providecommand{\atomsof}{\mathrm{atoms}}
\providecommand{\Red}{\rd{\mathsf{R}}}           % the red HF universe
   % red/black colour layer + FM symbols (providecommand-guarded)

\section{Future work: palettes and the generation of colour}
\label{sec:future}

\subsection{Arbitrarily many colours}
\label{sec:future:palette}

Two colours were a convenience, not a necessity. Replace the pair by any set $C$
of \emph{colours}; each $c\in C$ carries its own copy of the theory,
$\mathrm{Set}_c[X_c]$, with its own membership $\in_c$ and its own comprehension
$\{\cdots\}_c$. The only remaining datum is the \emph{nesting discipline}: which
colours may serve as atoms for which. This is a relation on $C$, and it is most
naturally presented as a directed graph --- equivalently a finite-state machine
--- with states $C$ and an edge $c\to c'$ whenever $c'$-sets are admissible
$c$-atoms. The atom supply of each colour is then the disjoint sum over its
outgoing edges,
\[
  X_c \;=\; \bigsqcup_{c\to c'}\ \mathrm{Set}_{c'}[X_{c'}],
\]
and the two-colour theory of this paper is the instance given by the two-state
automaton $\rd{\mathrm{red}} \rightleftarrows \bk{\mathrm{black}}$, with
$\rX=\bSetX$ and $\bX=\rSetX$.

Everything in the paper relativises to this picture. The well-formed cross-colour
nestings are the runs of the automaton; the local/global well-foundedness analysis
of Section~\ref{sec:nwf} reads off its structure, an acyclic colour graph giving a
globally well-founded universe and a cycle giving non-well-foundedness along that
cycle. (A self-loop $c\to c$ would let $c$-sets be their own atoms and collapse
the barrier; the two-cycle $\rd{\mathrm{red}}\rightleftarrows\bk{\mathrm{black}}$
is the smallest discipline that keeps every colour locally well-founded while
making the system globally non-well-founded.) This generalization is routine, and
we record it mainly to fix the slogan: \emph{a reflective set theory is an
automaton on colours}, and the constructions above are its two-state instance.

\subsection{Generating colours dynamically}
\label{sec:future:dynamic}

The genuinely interesting question we leave open. A fixed automaton fixes the
palette in advance. One would instead like colours to be \emph{generated on
demand}: a fresh colour spawned in the course of building a set, the palette
growing as construction proceeds, the automaton unfolding rather than given. The
difficulty is not in admitting infinitely many colours --- a countable palette is
unremarkable --- but in doing so \emph{elegantly}, with a generation discipline
internal to the theory, so that the mechanism producing colours is of a piece with
the mechanism producing sets and names.

There is a suggestive circularity to exploit, which is why we think the question
worth posing rather than dismissing. Generating a fresh colour is generating a
fresh name, in exactly the sense that
Sections~\ref{sec:fm}--\ref{sec:applications} reconstruct: a colour is drawn from
a supply, distinct from those already in play, and otherwise structureless. So a
reflective theory might be made to \emph{generate its own colours} by the very
nominal machinery it models --- the palette drawn, like the atoms, from a stream,
and the colour automaton itself becoming a nominal object in a theory one level
up. Whether this reflexive step can be taken without infinite regress, what its
fixed point looks like, and whether it lands back inside the spectrum of
Section~\ref{sec:nwf}, are the questions we find most worth pursuing. We do not
settle them here; we mark them as the natural sequel to this work.
.
%% rset.colors.tex
%% Red/black colour layer for the two intertwined copies, plus the FM symbols
%% used by rset.mainthm.tex and rset.applications.tex.
%%
%% EVERYTHING routes through \rd (red copy = player) and \bk (black copy =
%% opponent): redefine just those two macros to change the whole scheme (e.g.
%% to make the "black" copy dark blue for legibility).  Every definition is
%% \providecommand-guarded, so it is safe to \input this file from several
%% places and harmless if rset.local already defines these names.
%%
%% Requires \usepackage{xcolor} (already loaded in rset.tex).

% --- the two copies -------------------------------------------------------
\providecommand{\rd}[1]{\textcolor{red}{#1}}     % red copy  (player)
\providecommand{\bk}[1]{\textcolor{black}{#1}}   % black copy (opponent)

% --- theory names and atom supplies --------------------------------------
\providecommand{\rSetX}{\rd{\mathrm{Set}[X]}}    % red   Set[X]
\providecommand{\bSetX}{\bk{\mathrm{Set}[X]}}    % black Set[X]
\providecommand{\rX}{\rd{X}}                      % red   atom supply
\providecommand{\bX}{\bk{X}}                      % black atom supply

% --- empty set ------------------------------------------------------------
\providecommand{\remp}{\rd{\emptyset}}
\providecommand{\bemp}{\bk{\emptyset}}

% --- membership (relations) ----------------------------------------------
\providecommand{\rin}{\mathrel{\rd{\in}}}
\providecommand{\bin}{\mathrel{\bk{\in}}}
\providecommand{\rnin}{\mathrel{\rd{\notin}}}
\providecommand{\bnin}{\mathrel{\bk{\notin}}}

% --- binary operations ----------------------------------------------------
\providecommand{\rcup}{\mathbin{\rd{\cup}}}
\providecommand{\bcup}{\mathbin{\bk{\cup}}}
\providecommand{\rcap}{\mathbin{\rd{\cap}}}
\providecommand{\bcap}{\mathbin{\bk{\cap}}}
\providecommand{\rsm}{\mathbin{\rd{\setminus}}}
\providecommand{\bsm}{\mathbin{\bk{\setminus}}}

% --- comprehension braces of each theory ---------------------------------
\providecommand{\rbr}[1]{\rd{\{}\,#1\,\rd{\}}}
\providecommand{\bbr}[1]{\bk{\{}\,#1\,\bk{\}}}

% --- colour-blind membership (sees through the colour barrier) ------------
\providecommand{\cbin}{\mathrel{\in_{\!\star}}}

% --- FM symbols (shared with rset.stdfmset.tex) --------------------------
\providecommand{\Atoms}{\mathbb{A}}
\providecommand{\Sym}{\mathrm{Sym}}
\providecommand{\pow}{\mathcal{P}}
\providecommand{\supp}{\mathrm{supp}}
\providecommand{\fs}{\mathrm{fs}}
\providecommand{\new}{\mathbin{\reflectbox{\ensuremath{\mathsf{N}}}}}
\providecommand{\atomsof}{\mathrm{atoms}}
\providecommand{\Red}{\rd{\mathsf{R}}}           % the red HF universe
   % red/black colour layer + FM symbols (providecommand-guarded)

\section{Future work: palettes and the generation of colour}
\label{sec:future}

\subsection{Arbitrarily many colours}
\label{sec:future:palette}

Two colours were a convenience, not a necessity. Replace the pair by any set $C$
of \emph{colours}; each $c\in C$ carries its own copy of the theory,
$\mathrm{Set}_c[X_c]$, with its own membership $\in_c$ and its own comprehension
$\{\cdots\}_c$. The only remaining datum is the \emph{nesting discipline}: which
colours may serve as atoms for which. This is a relation on $C$, and it is most
naturally presented as a directed graph --- equivalently a finite-state machine
--- with states $C$ and an edge $c\to c'$ whenever $c'$-sets are admissible
$c$-atoms. The atom supply of each colour is then the disjoint sum over its
outgoing edges,
\[
  X_c \;=\; \bigsqcup_{c\to c'}\ \mathrm{Set}_{c'}[X_{c'}],
\]
and the two-colour theory of this paper is the instance given by the two-state
automaton $\rd{\mathrm{red}} \rightleftarrows \bk{\mathrm{black}}$, with
$\rX=\bSetX$ and $\bX=\rSetX$.

Everything in the paper relativises to this picture. The well-formed cross-colour
nestings are the runs of the automaton; the local/global well-foundedness analysis
of Section~\ref{sec:nwf} reads off its structure, an acyclic colour graph giving a
globally well-founded universe and a cycle giving non-well-foundedness along that
cycle. (A self-loop $c\to c$ would let $c$-sets be their own atoms and collapse
the barrier; the two-cycle $\rd{\mathrm{red}}\rightleftarrows\bk{\mathrm{black}}$
is the smallest discipline that keeps every colour locally well-founded while
making the system globally non-well-founded.) This generalization is routine, and
we record it mainly to fix the slogan: \emph{a reflective set theory is an
automaton on colours}, and the constructions above are its two-state instance.

\subsection{Generating colours dynamically}
\label{sec:future:dynamic}

The genuinely interesting question we leave open. A fixed automaton fixes the
palette in advance. One would instead like colours to be \emph{generated on
demand}: a fresh colour spawned in the course of building a set, the palette
growing as construction proceeds, the automaton unfolding rather than given. The
difficulty is not in admitting infinitely many colours --- a countable palette is
unremarkable --- but in doing so \emph{elegantly}, with a generation discipline
internal to the theory, so that the mechanism producing colours is of a piece with
the mechanism producing sets and names.

There is a suggestive circularity to exploit, which is why we think the question
worth posing rather than dismissing. Generating a fresh colour is generating a
fresh name, in exactly the sense that
Sections~\ref{sec:fm}--\ref{sec:applications} reconstruct: a colour is drawn from
a supply, distinct from those already in play, and otherwise structureless. So a
reflective theory might be made to \emph{generate its own colours} by the very
nominal machinery it models --- the palette drawn, like the atoms, from a stream,
and the colour automaton itself becoming a nominal object in a theory one level
up. Whether this reflexive step can be taken without infinite regress, what its
fixed point looks like, and whether it lands back inside the spectrum of
Section~\ref{sec:nwf}, are the questions we find most worth pursuing. We do not
settle them here; we mark them as the natural sequel to this work.
 from rset.tex.
%% rset.colors.tex
%% Red/black colour layer for the two intertwined copies, plus the FM symbols
%% used by rset.mainthm.tex and rset.applications.tex.
%%
%% EVERYTHING routes through \rd (red copy = player) and \bk (black copy =
%% opponent): redefine just those two macros to change the whole scheme (e.g.
%% to make the "black" copy dark blue for legibility).  Every definition is
%% \providecommand-guarded, so it is safe to \input this file from several
%% places and harmless if rset.local already defines these names.
%%
%% Requires \usepackage{xcolor} (already loaded in rset.tex).

% --- the two copies -------------------------------------------------------
\providecommand{\rd}[1]{\textcolor{red}{#1}}     % red copy  (player)
\providecommand{\bk}[1]{\textcolor{black}{#1}}   % black copy (opponent)

% --- theory names and atom supplies --------------------------------------
\providecommand{\rSetX}{\rd{\mathrm{Set}[X]}}    % red   Set[X]
\providecommand{\bSetX}{\bk{\mathrm{Set}[X]}}    % black Set[X]
\providecommand{\rX}{\rd{X}}                      % red   atom supply
\providecommand{\bX}{\bk{X}}                      % black atom supply

% --- empty set ------------------------------------------------------------
\providecommand{\remp}{\rd{\emptyset}}
\providecommand{\bemp}{\bk{\emptyset}}

% --- membership (relations) ----------------------------------------------
\providecommand{\rin}{\mathrel{\rd{\in}}}
\providecommand{\bin}{\mathrel{\bk{\in}}}
\providecommand{\rnin}{\mathrel{\rd{\notin}}}
\providecommand{\bnin}{\mathrel{\bk{\notin}}}

% --- binary operations ----------------------------------------------------
\providecommand{\rcup}{\mathbin{\rd{\cup}}}
\providecommand{\bcup}{\mathbin{\bk{\cup}}}
\providecommand{\rcap}{\mathbin{\rd{\cap}}}
\providecommand{\bcap}{\mathbin{\bk{\cap}}}
\providecommand{\rsm}{\mathbin{\rd{\setminus}}}
\providecommand{\bsm}{\mathbin{\bk{\setminus}}}

% --- comprehension braces of each theory ---------------------------------
\providecommand{\rbr}[1]{\rd{\{}\,#1\,\rd{\}}}
\providecommand{\bbr}[1]{\bk{\{}\,#1\,\bk{\}}}

% --- colour-blind membership (sees through the colour barrier) ------------
\providecommand{\cbin}{\mathrel{\in_{\!\star}}}

% --- FM symbols (shared with rset.stdfmset.tex) --------------------------
\providecommand{\Atoms}{\mathbb{A}}
\providecommand{\Sym}{\mathrm{Sym}}
\providecommand{\pow}{\mathcal{P}}
\providecommand{\supp}{\mathrm{supp}}
\providecommand{\fs}{\mathrm{fs}}
\providecommand{\new}{\mathbin{\reflectbox{\ensuremath{\mathsf{N}}}}}
\providecommand{\atomsof}{\mathrm{atoms}}
\providecommand{\Red}{\rd{\mathsf{R}}}           % the red HF universe
   % red/black colour layer + FM symbols (providecommand-guarded)

\section{Conclusions and future work}
\label{sec:conc}
\label{sec:future}

We set out to reconcile two things usually held apart: a set theory with atoms,
of the kind nominal techniques require, and the austerity of building everything
from the empty set. The reconciliation is the knot. Taking two copies of set
theory --- a {\rd red} and a {\bk black} --- and tying them so that each colour's
atoms are the other colour's sets, we obtain an infinite supply of atoms without
positing a single one: every atom is a set of the dual colour, manufactured from
$\varnothing$, and the infinite risk that Chaitin's accounting charges to a
primitive infinity of atoms is refinanced as the ordinary risk of one empty set
and one successor.

Around this knot we have assembled a connected picture. The formal theory above
fixes the two copies and ties them; a self-contained sketch of
Fraenkel--Mostowski set theory (Section~\ref{sec:fm}) makes the target precise;
and the representation theorem (Section~\ref{sec:fmmodel}) exhibits the red
hereditarily finite sets over the black-set atoms as a model of the nominal
hereditarily finite sets, in which finite support --- imposed as an axiom in
$\mathsf{FM}$ --- is instead a theorem. The boundary between the finitary core and
its infinitary extension proved to be the exact ledger line between what the
construction \emph{derives} --- atoms, foundation, finite support --- and what it
must still \emph{assume}: a single $\omega$, and the finite-support cut, once
actual infinities are admitted.

The construction earns its keep in three further ways. It founds the
Pitts--Gabbay nominal techniques (Section~\ref{sec:applications}), and on exactly
the hereditarily finite fragment that the realised nominal languages inhabit, so
the restriction that bounds the set theory is the natural home of the computing;
it suggests, through reflective names, a more expressive theory of names in which
the nominal discipline reappears as an equivariant fragment; and it serves as a
bridge (Section~\ref{sec:nwf}) between the well-founded theories $\mathsf{ZF}$ and
$\mathsf{FM}$ and Aczel's non-well-founded universe, by making
non-well-foundedness not an axiom but an emergent, global property of two locally
well-founded theories. The road from this programmatic account to a rigorous one
runs, we have argued, through algebraic set theory (Section~\ref{sec:rigor}),
whose initial-$\mathsf{ZF}$-algebra theorem is purpose-built to supply the knotted
fixpoint and, with it, the membership and extensionality structure the
construction needs.

Several directions remain open; we record the ones we find most pressing.

\subsection{Arbitrarily many colours}
\label{sec:future:palette}

Two colours were a convenience, not a necessity. Replace the pair by any set $C$
of \emph{colours}; each $c\in C$ carries its own copy of the theory,
$\mathrm{Set}_c[X_c]$, with its own membership $\in_c$ and its own comprehension
$\{\cdots\}_c$. The only remaining datum is the \emph{nesting discipline}: which
colours may serve as atoms for which. This is a relation on $C$, and it is most
naturally presented as a directed graph --- equivalently a finite-state machine
--- with states $C$ and an edge $c\to c'$ whenever $c'$-sets are admissible
$c$-atoms. The atom supply of each colour is then the disjoint sum over its
outgoing edges,
\[
  X_c \;=\; \bigsqcup_{c\to c'}\ \mathrm{Set}_{c'}[X_{c'}],
\]
and the two-colour theory of this paper is the instance given by the two-state
automaton $\rd{\mathrm{red}} \rightleftarrows \bk{\mathrm{black}}$, with
$\rX=\bSetX$ and $\bX=\rSetX$.

Everything in the paper relativises to this picture. The well-formed cross-colour
nestings are the runs of the automaton; the local/global well-foundedness analysis
of Section~\ref{sec:nwf} reads off its structure, an acyclic colour graph giving a
globally well-founded universe and a cycle giving non-well-foundedness along that
cycle. (A self-loop $c\to c$ would let $c$-sets be their own atoms and collapse
the barrier; the two-cycle $\rd{\mathrm{red}}\rightleftarrows\bk{\mathrm{black}}$
is the smallest discipline that keeps every colour locally well-founded while
making the system globally non-well-founded.) This generalization is routine, and
we record it mainly to fix the slogan: \emph{a reflective set theory is an
automaton on colours}, and the construction above is its two-state instance.

\subsection{Generating colours dynamically}
\label{sec:future:dynamic}

The genuinely interesting question we leave open. A fixed automaton fixes the
palette in advance. One would instead like colours to be \emph{generated on
demand}: a fresh colour spawned in the course of building a set, the palette
growing as construction proceeds, the automaton unfolding rather than given. The
difficulty is not in admitting infinitely many colours --- a countable palette is
unremarkable --- but in doing so \emph{elegantly}, with a generation discipline
internal to the theory, so that the mechanism producing colours is of a piece with
the mechanism producing sets and names.

There is a suggestive circularity to exploit, which is why we think the question
worth posing rather than dismissing. Generating a fresh colour is generating a
fresh name, in exactly the sense that
Sections~\ref{sec:fm}--\ref{sec:applications} reconstruct: a colour is drawn from
a supply, distinct from those already in play, and otherwise structureless. So a
reflective theory might be made to \emph{generate its own colours} by the very
nominal machinery it models --- the palette drawn, like the atoms, from a stream,
and the colour automaton itself becoming a nominal object in a theory one level
up. Whether this reflexive step can be taken without infinite regress, what its
fixed point looks like, and whether it lands back inside the spectrum of
Section~\ref{sec:nwf}, are the questions we find most worth pursuing. We do not
settle them here; we mark them as the natural sequel to this work.

\subsection{A semantics for the rho calculus}
\label{sec:future:rho}

We close with a conjecture rather than a construction. Aczel gave a semantics for
Milner's $\mathsf{CCS}$ inside non-well-founded set theory \cite{aczel1988nwf}: a
process is identified with the set of its labelled transitions,
$\{\,(a,P') : P \xrightarrow{a} P'\,\}$, so that a looping process is a
non-well-founded set, bisimilar processes are equal, and the anti-foundation
axiom's solution lemma supplies and individuates processes as the unique
solutions of their behavioural equations. We surmise that the knotted universe
developed here is the natural home for the same construction carried out for the
reflective higher-order calculus, the rho calculus
\cite{DBLP:journals/entcs/MeredithR05}.

The fit is suggestive on two counts. First, the non-well-foundedness that Aczel
imports by axiom is here available for free, as the global, colour-blind structure
of the knotted universe (Section~\ref{sec:nwf}); the recursion that process
behaviour requires is the recursion the knot already carries. Second --- and this
is what sets rho apart from $\mathsf{CCS}$ and the $\pi$-calculus --- rho's names
are not primitive but \emph{reflected}: the quote operator $@$ turns a process
into a name, the dereference operator $*$ turns a name back into the process it
quotes, and $*@P \equiv P$. In the knotted universe a name may be taken to be an
\emph{atom}, and atoms here are constructed sets of the dual colour, opaque until
one crosses the colour barrier (Sections~\ref{sec:fmmodel:atoms}
and~\ref{sec:applications:reflective}). Reading $@$ as \emph{atom construction}
--- a process, presented as a set, is sealed into an atom of the opposite colour
--- and $*$ as \emph{atom unpacking} --- the barrier is crossed to recover the set
the atom seals --- the reflection law $*@P \equiv P$ becomes the round trip of
these two operations.

We do not carry the construction out, and we do not claim its details are routine:
communication, the substitution of names for names, and rho's equational theory
would each have to be matched against their set-theoretic shadows, and the
well-founded versus non-well-founded reading (Section~\ref{sec:rigor:fixpoint})
settled. We claim only that the ingredients line up --- a calculus whose names are
quoted processes asks for a set theory whose atoms are sets, and whose dynamics
asks for non-well-foundedness, and the knotted universe offers both at once --- and
we conjecture that such a semantics exists and is more convenient than one built
over an unstructured atom supply, precisely because the $@$/$*$ reflection it must
model is already the construction's defining move.
